\documentclass[11pt,a4paper]{article}

% --- Packages ---
\usepackage[utf8]{inputenc}
\usepackage[T1]{fontenc}
\usepackage{amsmath,amssymb,amsthm}
\usepackage{physics}
\usepackage{geometry}
\usepackage{enumitem}
\usepackage{hyperref}
\usepackage{fancyhdr}
\usepackage{mathtools}

\geometry{margin=1in}
\pagestyle{fancy}
\fancyhf{}
\rhead{Quantum Spin Lecture Notes}
\lhead{Lecture 5: Schr\"odinger's Equation \& Matrix Exponentiation}
\cfoot{\thepage}

% --- Theorem Environments ---
\theoremstyle{definition}
\newtheorem{definition}{Definition}[section]
\newtheorem{example}{Example}[section]
\theoremstyle{plain}
\newtheorem{theorem}{Theorem}[section]
\newtheorem{proposition}{Proposition}[section]
\newtheorem{corollary}{Corollary}[section]
\theoremstyle{remark}
\newtheorem{remark}{Remark}[section]

\title{\textbf{Quantum Spin (5): Schr\"odinger's Equation \\ \& Matrix Exponentiation}}
\author{Lecture Notes on Quantum Spin}
\date{}

\begin{document}
\maketitle
\tableofcontents
\newpage

% ============================================================
\section{Setup and Review}
% ============================================================

A point-like electron with spin state $\ket{\psi(t)} = \begin{pmatrix}c_1(t)\\c_2(t)\end{pmatrix}$ sits in an external magnetic field $\vec{B}$. The Schr\"odinger equation governing time evolution is:
\[
    i\hbar\frac{d}{dt}\ket{\psi(t)} = H\ket{\psi(t)},
\]
equivalently:
\[
    \frac{d}{dt}\ket{\psi(t)} = -\frac{i}{\hbar}H\ket{\psi(t)}.
\]

The Hamiltonian is $H = \mathcal{E}\,(\hat{n}\cdot\vec{\sigma})$, where:
\begin{itemize}
    \item $\vec{B} = B\hat{n}$ ($\hat{n}$ is a unit vector, $B = |\vec{B}|$),
    \item $\mathcal{E} = \frac{ge\hbar B}{4m}$ is a constant with dimensions of energy,
    \item $\hat{n}\cdot\vec{\sigma} = n_x\sigma_x + n_y\sigma_y + n_z\sigma_z$ is a $2\times 2$ matrix.
\end{itemize}

% ============================================================
\section{Method 1: Direct Solution}
% ============================================================

\subsection{Example: $\vec{B}$ Along the $y$-Axis}

Let $\hat{n} = \hat{y}$, so $H = \mathcal{E}\,\sigma_y = \begin{pmatrix}0 & -i\mathcal{E} \\ i\mathcal{E} & 0\end{pmatrix}$.

Initial condition: $\ket{\psi(0)} = \begin{pmatrix}1\\0\end{pmatrix}$ (spin up along $z$).

The Schr\"odinger equation becomes two coupled ODEs:
\begin{align}
    \dot{c}_1(t) &= -\frac{\mathcal{E}}{\hbar}\,c_2(t), \label{eq:ode1}\\[4pt]
    \dot{c}_2(t) &= +\frac{\mathcal{E}}{\hbar}\,c_1(t). \label{eq:ode2}
\end{align}

\subsection{Decoupling}

Differentiate \eqref{eq:ode1} and substitute \eqref{eq:ode2}:
\[
    \ddot{c}_1 = -\frac{\mathcal{E}}{\hbar}\dot{c}_2 = -\frac{\mathcal{E}^2}{\hbar^2}\,c_1.
\]
Similarly, $\ddot{c}_2 = -\frac{\mathcal{E}^2}{\hbar^2}\,c_2$. These are simple harmonic oscillator equations with angular frequency $\omega = \mathcal{E}/\hbar$.

\subsection{General Solution and Initial Conditions}

The general solutions are:
\begin{align}
    c_1(t) &= A\cos\!\left(\frac{\mathcal{E}}{\hbar}t\right) + B\sin\!\left(\frac{\mathcal{E}}{\hbar}t\right), \\
    c_2(t) &= C\cos\!\left(\frac{\mathcal{E}}{\hbar}t\right) + D\sin\!\left(\frac{\mathcal{E}}{\hbar}t\right).
\end{align}

Using initial conditions $c_1(0) = 1$, $c_2(0) = 0$ and the original ODEs (guess and check):
\[
    \boxed{\ket{\psi(t)} = \begin{pmatrix}\cos\!\left(\dfrac{\mathcal{E}}{\hbar}t\right) \\[8pt] \sin\!\left(\dfrac{\mathcal{E}}{\hbar}t\right)\end{pmatrix}.}
\]

\subsection{Verification}

\begin{itemize}
    \item $|c_1|^2 + |c_2|^2 = \cos^2(\mathcal{E}t/\hbar) + \sin^2(\mathcal{E}t/\hbar) = 1$ for all $t$. $\checkmark$
    \item $\dot{c}_1 = -(\mathcal{E}/\hbar)\sin(\mathcal{E}t/\hbar) = -(\mathcal{E}/\hbar)c_2$. $\checkmark$
    \item $\dot{c}_2 = +(\mathcal{E}/\hbar)\cos(\mathcal{E}t/\hbar) = +(\mathcal{E}/\hbar)c_1$. $\checkmark$
\end{itemize}

\subsection{Physical Interpretation: Larmor Precession}

Evaluating $\ket{\psi(t)}$ at key times (with $T_0 = \hbar/\mathcal{E}$):

\begin{center}
\renewcommand{\arraystretch}{1.5}
\begin{tabular}{c|c|c}
    $t$ & $\ket{\psi(t)}$ & Direction \\ \hline
    $0$ & $(1,\; 0)^\top$ & $+z$ \\
    $\frac{\pi}{4}T_0$ & $\frac{1}{\sqrt{2}}(1,\; 1)^\top$ & $+x$ \\
    $\frac{\pi}{2}T_0$ & $(0,\; 1)^\top$ & $-z$ \\
    $\frac{3\pi}{4}T_0$ & $\frac{1}{\sqrt{2}}(-1,\; 1)^\top$ & $-x$ (up to phase) \\
    $\pi T_0$ & $(-1,\; 0)^\top$ & $+z$ (negated)
\end{tabular}
\end{center}

The Bloch vector traces a circle in the $xz$-plane, precessing around the $y$-axis (the direction of $\vec{B}$). After a full $2\pi$ rotation, $\ket{\psi} \to -\ket{\psi}$ --- the state is \emph{negated} (a signature of spin-$\frac{1}{2}$).

% ============================================================
\section{Method 2: Eigendecomposition}
% ============================================================

\subsection{Finding Eigenstates of $H$}

Since $H = \mathcal{E}\,\sigma_y$, the eigenstates of $H$ are those of $\sigma_y$:
\begin{align}
    H\ket{\uparrow_y} &= +\mathcal{E}\ket{\uparrow_y}, \\
    H\ket{\downarrow_y} &= -\mathcal{E}\ket{\downarrow_y}.
\end{align}

\subsection{Eigenstates Yield Immediate Solutions}

If $H\ket{\phi} = E\ket{\phi}$, then:
\[
    \ket{\phi(t)} = e^{-iEt/\hbar}\ket{\phi}
\]
solves the Schr\"odinger equation. (Verification: $i\hbar\frac{d}{dt}\ket{\phi(t)} = E\,e^{-iEt/\hbar}\ket{\phi} = H\ket{\phi(t)}$.)

The two solutions are:
\begin{align}
    \ket{\psi_1(t)} &= e^{-i\mathcal{E}t/\hbar}\ket{\uparrow_y}, \\
    \ket{\psi_2(t)} &= e^{+i\mathcal{E}t/\hbar}\ket{\downarrow_y}.
\end{align}

\subsection{Linearity and Superposition}

Because the Schr\"odinger equation is \textbf{linear}, any linear combination of solutions is also a solution:
\[
    \ket{\psi(t)} = \alpha\ket{\psi_1(t)} + \beta\ket{\psi_2(t)}.
\]

To match the initial condition $\ket{\psi(0)} = \begin{pmatrix}1\\0\end{pmatrix}$:
\[
    \begin{pmatrix}1\\0\end{pmatrix} = \alpha\ket{\uparrow_y} + \beta\ket{\downarrow_y} \implies \alpha = \beta = \frac{1}{\sqrt{2}}.
\]

Therefore:
\[
    \ket{\psi(t)} = \frac{1}{\sqrt{2}}e^{-i\mathcal{E}t/\hbar}\ket{\uparrow_y} + \frac{1}{\sqrt{2}}e^{+i\mathcal{E}t/\hbar}\ket{\downarrow_y}.
\]

Expanding and using Euler's formula yields the same answer as Method 1:
\[
    \ket{\psi(t)} = \begin{pmatrix}\cos(\mathcal{E}t/\hbar) \\ \sin(\mathcal{E}t/\hbar)\end{pmatrix}. \quad\checkmark
\]

% ============================================================
\section{Method 3: Matrix Exponentiation}
% ============================================================

\subsection{Motivation via Time Stepping}

From a first-order Taylor expansion:
\[
    \ket{\psi(t + \delta t)} \approx \ket{\psi(t)} + \delta t\,\frac{d}{dt}\ket{\psi(t)} = \left(I - \frac{i}{\hbar}\delta t\, H\right)\ket{\psi(t)}.
\]

Breaking time $t$ into $n$ small steps of size $t/n$:
\[
    \ket{\psi(t)} \approx \left(I - \frac{i}{\hbar}\frac{t}{n}H\right)^n \ket{\psi(0)}.
\]

Taking $n \to \infty$ and using the identity $e^A = \lim_{n\to\infty}\left(I + \frac{A}{n}\right)^n$:
\[
    \boxed{\ket{\psi(t)} = e^{-iHt/\hbar}\ket{\psi(0)}.}
\]

\subsection{Rigorous Definition of the Matrix Exponential}

\begin{definition}[Matrix Exponential]
    For a square matrix $A$:
    \[
        e^A \coloneqq \sum_{n=0}^{\infty} \frac{A^n}{n!} = I + A + \frac{A^2}{2!} + \frac{A^3}{3!} + \cdots
    \]
    where $A^0 = I$ (the identity matrix).
\end{definition}

\subsection{Key Property}

\begin{theorem}[Differentiation of the Matrix Exponential]
    \[
        \frac{d}{dt}e^{tA} = A\,e^{tA}.
    \]
\end{theorem}

\begin{proof}
    \[
        \frac{d}{dt}e^{tA} = \frac{d}{dt}\sum_{n=0}^{\infty}\frac{t^n A^n}{n!} = \sum_{n=1}^{\infty}\frac{n\,t^{n-1}A^n}{n!} = A\sum_{n=1}^{\infty}\frac{t^{n-1}A^{n-1}}{(n-1)!} = A\,e^{tA}. \qedhere
    \]
\end{proof}

This immediately implies $\ket{\psi(t)} = e^{-iHt/\hbar}\ket{\psi(0)}$ solves the Schr\"odinger equation.

\subsection{Computing Matrix Exponentials}

\subsubsection{Method A: Diagonalization}

If $A = M D M^{-1}$ where $D$ is diagonal, then $A^n = MD^nM^{-1}$ and:
\[
    e^A = M\,e^D\,M^{-1}, \quad \text{where } e^D = \operatorname{diag}(e^{\lambda_1}, e^{\lambda_2}, \ldots).
\]

\subsubsection{Method B: Using $A^2 = -I$ (Generalized Euler Formula)}

If $A$ is a matrix satisfying $A^2 = -I$, then:
\[
    \boxed{e^{\theta A} = \cos\theta\, I + \sin\theta\, A.}
\]

\begin{proof}
    Separate the Taylor series into even and odd powers:
    \[
        e^{\theta A} = \sum_{n=0}^{\infty}\frac{(\theta A)^n}{n!} = \underbrace{\sum_{k=0}^{\infty}\frac{(-1)^k\theta^{2k}}{(2k)!}}_{=\cos\theta}\,I + \underbrace{\sum_{k=0}^{\infty}\frac{(-1)^k\theta^{2k+1}}{(2k+1)!}}_{=\sin\theta}\,A,
    \]
    using $A^{2k} = (-I)^k = (-1)^kI$ and $A^{2k+1} = (-1)^kA$. \qedhere
\end{proof}

\begin{remark}
    This is the matrix analog of Euler's formula $e^{i\theta} = \cos\theta + i\sin\theta$, since $i^2 = -1$ parallels $A^2 = -I$.
\end{remark}

% ============================================================
\section{Applying Matrix Exponentiation to Spin}
% ============================================================

\subsection{The Time Evolution Operator}

We need to compute $e^{-iHt/\hbar}$ where $H = \mathcal{E}(\hat{n}\cdot\vec{\sigma})$.

Define $A = -i(\hat{n}\cdot\vec{\sigma})$. Then $-iHt/\hbar = (\mathcal{E}t/\hbar)\,A$.

\subsection{Verifying $A^2 = -I$}

\begin{align}
    A^2 &= (-i)^2(\hat{n}\cdot\vec{\sigma})^2 = -(\hat{n}\cdot\vec{\sigma})^2.
\end{align}

Computing $(\hat{n}\cdot\vec{\sigma})^2$ explicitly (with $|\hat{n}|=1$):
\[
    (\hat{n}\cdot\vec{\sigma})^2 = \begin{pmatrix}n_z & n_x-in_y \\ n_x+in_y & -n_z\end{pmatrix}^2 = \begin{pmatrix}n_x^2+n_y^2+n_z^2 & 0 \\ 0 & n_x^2+n_y^2+n_z^2\end{pmatrix} = I.
\]

Therefore $A^2 = -I$. $\checkmark$

\subsection{The General Formula}

Applying the generalized Euler formula:
\[
    \boxed{e^{-iHt/\hbar} = \cos\!\left(\frac{\mathcal{E}t}{\hbar}\right)I + \sin\!\left(\frac{\mathcal{E}t}{\hbar}\right)\bigl(-i\,\hat{n}\cdot\vec{\sigma}\bigr).}
\]

\subsection{Verification with the $y$-Axis Example}

For $\hat{n} = \hat{y}$: $-i\sigma_y = -i\begin{pmatrix}0&-i\\i&0\end{pmatrix} = \begin{pmatrix}0&-1\\1&0\end{pmatrix}$.

\[
    e^{-iHt/\hbar} = \begin{pmatrix}\cos(\mathcal{E}t/\hbar) & -\sin(\mathcal{E}t/\hbar) \\ \sin(\mathcal{E}t/\hbar) & \cos(\mathcal{E}t/\hbar)\end{pmatrix}.
\]

Acting on $\ket{\psi(0)} = \begin{pmatrix}1\\0\end{pmatrix}$:
\[
    \ket{\psi(t)} = \begin{pmatrix}\cos(\mathcal{E}t/\hbar)\\\sin(\mathcal{E}t/\hbar)\end{pmatrix}. \quad\checkmark
\]

% ============================================================
\section{Time Evolution as Rotation}
% ============================================================

\subsection{$\vec{B}$ Along the $z$-Axis}

If $\hat{n} = \hat{z}$, then $H = \mathcal{E}\,\sigma_z$ and:
\[
    e^{-iHt/\hbar} = \begin{pmatrix}e^{-i\mathcal{E}t/\hbar} & 0 \\ 0 & e^{+i\mathcal{E}t/\hbar}\end{pmatrix}.
\]

Acting on a general state $\ket{\psi} = \begin{pmatrix}\cos\frac{\theta}{2} \\ e^{i\phi}\sin\frac{\theta}{2}\end{pmatrix}$:
\[
    \ket{\psi(t)} = \begin{pmatrix}e^{-i\mathcal{E}t/\hbar}\cos\frac{\theta}{2} \\ e^{i(\phi + 2\mathcal{E}t/\hbar)}\sin\frac{\theta}{2}\end{pmatrix} \sim \begin{pmatrix}\cos\frac{\theta}{2} \\ e^{i(\phi + 2\mathcal{E}t/\hbar)}\sin\frac{\theta}{2}\end{pmatrix},
\]
where $\sim$ denotes equality up to a global phase. The polar angle $\theta$ is unchanged; the azimuthal angle increases as $\phi \to \phi + 2\mathcal{E}t/\hbar$. The Bloch vector \textbf{rotates around the $z$-axis}.

\subsection{General Rotation Formula}

\begin{theorem}[Spin Rotation]
    The matrix
    \[
        U(\hat{n}, \varphi) = e^{-i\varphi\,\hat{n}\cdot\vec{\sigma}/2}
    \]
    implements a \textbf{rotation of the Bloch vector} around the axis $\hat{n}$ by angle $\varphi$.
\end{theorem}

\subsection{Two Key Takeaways}

\begin{enumerate}
    \item \textbf{Time evolution in a magnetic field \emph{is} a rotation.} The Bloch vector precesses around $\hat{n}$ (the direction of $\vec{B}$) at a rate determined by $\mathcal{E}/\hbar$.
    \item \textbf{The matrix exponential gives an explicit rotation formula.} Once $e^{-iHt/\hbar}$ is computed, applying it to \emph{any} initial state immediately yields the time-evolved state.
\end{enumerate}

\subsection{Connection to Quantum Computing}

In the language of quantum computation:
\begin{itemize}
    \item The time evolution operator $U(t) = e^{-iHt/\hbar}$ is a \textbf{quantum gate}.
    \item A spin state is a \textbf{qubit}.
    \item Any computational operation on a qubit can be realized by choosing an appropriate magnetic field direction $\hat{n}$ and duration $t$.
\end{itemize}

\end{document}
